% Book-local P07.1 theorem, 13 September 2026.
% This theorem combines the locally constant Stone-parameter mechanism with
% the finite ambient-star x convex-polyhedron mechanism.  It makes no
% novelty or priority claim.
\begin{theorem}
\label{thm:p071-stone-finite-star-polyhedral-collapse}
Let \(Z\) be a locally compact Hausdorff space with a basis of clopen
subsets.  Let \(E\) be a finite-dimensional real vector space and let
\[
 S=\{0\}\cup\coprod_{i=1}^{m}(0,r_i)v_i\subset E,
 \qquad 1\leq m<\infty,
\]
where \(r_i>0\), \(v_i\ne0\), and the positive rays
\(\mathbf R_{>0}v_i\) are pairwise distinct.  Let \(P\) be a nonempty
convex polyhedron in a finite-dimensional real affine space, put
\(A=\operatorname{Aff}(P)\), and set
\[
                              X=Z\times S\times P.
\]

For an ordinary open \(W\subset X\), let \(\mathcal G_{X,0}(W)\) be the
real-valued functions with the following local description.  After
shrinking the \(Z\)-variable to a clopen neighbourhood, the function is
independent of that variable and is the restriction of a smooth ambient
function on \(E\times A\) whose restriction
\[
                 (t,y)\longmapsto g(tv_i,y)
\]
has positive-semidefinite full Hessian on every branch chamber on which it
is defined.  Let \(\mathcal P_{X,1}\) be the sheaf of local pointwise
decreasing limits of sequences in \(\mathcal G_{X,0}\), and let
\(\mathcal P_X\) be the smallest subsheaf of
\(\mathbf R\cup\{-\infty\}\)-valued functions which contains
\(\mathcal G_{X,0}\) and is closed under such limits.

Define \(\mathcal C_X\) by the following local condition.  A section \(u\)
belongs to \(\mathcal C_X\) near \((z_0,0,y_0)\) if, after shrinking to a
product of a compact clopen neighbourhood \(C\ni z_0\), a smaller finite
substar \(S'\ni0\), and a relatively open convex neighbourhood
\(V\ni y_0\) in \(P\), there are a compact metrizable zero-dimensional
space \(Y\), a continuous surjection \(q:C\twoheadrightarrow Y\), and an
upper semicontinuous function
\[
        \bar u:Y\times S'\times V
             \longrightarrow\mathbf R\cup\{-\infty\}
\]
such that
\begin{equation}
 u(z,s,y)=\bar u(q(z),s,y),
 \label{eq:p071-stone-star-factorization}
\end{equation}
and, for every \(\zeta\in Y\), the function
\(\bar u_\zeta(s,y):=\bar u(\zeta,s,y)\) satisfies:
\begin{enumerate}
\item on each branch, with the common value at \(t=0\), the function
  \((t,y)\mapsto\bar u_\zeta(tv_i,y)\) is extended-real jointly convex on
  every convex subset of its domain;
\item for every finite family of points
  \((t_av_{i(a)},y_a)\) and weights \(\alpha_a\geq0\) whose joining
  segments from \((0,y_*)\) remain in the domain and which satisfy
  \[
    \sum_a\alpha_a=1,
    \qquad \sum_a\alpha_at_av_{i(a)}=0,
    \qquad \sum_a\alpha_ay_a=y_*,
  \]
  one has the product junction inequality
  \begin{equation}
     \bar u_\zeta(0,y_*)
        \leq\sum_a\alpha_a
                    \bar u_\zeta(t_av_{i(a)},y_a).
   \label{eq:p071-stone-star-junction}
  \end{equation}
  Terms with zero coefficient are omitted from the extended-real sum.
\end{enumerate}
Away from the node, the same definition is used after restricting to the
single branch containing the point; condition
\eqref{eq:p071-stone-star-junction} is then vacuous.

Then \(\mathcal C_X\) is a sheaf and
\begin{equation}
                 \mathcal P_{X,1}=\mathcal P_X=\mathcal C_X.
 \label{eq:p071-stone-star-collapse}
\end{equation}
More precisely, every germ in these sheaves has a pointwise decreasing
approximation by members of \(\mathcal G_{X,0}\) on one fixed smaller
product neighbourhood.  At a node point \((z_0,0,y_0)\), the neighbourhood
may be chosen as \(C\times S''\times V'\), where \(C\) is compact clopen,
\(S''\) is a finite substar with one positive collar on every branch, and
\(V'\) is a relatively open polyhedral neighbourhood in \(P\).

The theorem permits simultaneously a nonmetrizable Stone parameter,
genuine finite branching, nonconstant affine directions, and all faces and
corners of one convex polyhedron.  It concerns the displayed product
generator sheaf.  It makes no assertion for nonconstant affine directions
on the Stone factor, an infinite non-CLD star, separately convex
generators, a nonambient branch structure, or a nonconvex chamber gluing.
Consequently it narrows but does not retire P07.1.
\end{theorem}
